% Pacotes básicos
% ----------------------------------------------------------
%\usepackage{lmodern}			% Usa a fonte Latin Modern
\usepackage[T1]{fontenc}		% Selecao de codigos de fonte.
\usepackage[utf8]{inputenc}		% Codificacao do documento (conversão automática dos acentos)
\usepackage{lastpage}			% Usado pela Ficha catalográfica
\usepackage{indentfirst}		% Indenta o primeiro parágrafo de cada seção.
\usepackage{color}		    	% Controle das cores
\usepackage{graphicx}			% Inclusão de gráficos
\usepackage{microtype} 			% para melhorias de justificação
\usepackage{ifthen}		    	% para montar condicionais
\PassOptionsToPackage{english,brazil}{babel}
\usepackage{babel} % para utilizar termos em portugues
\usepackage[final]{pdfpages}    % para incluir páginas de arquivos pdf
\usepackage{lipsum}				% para geração de dummy text
\usepackage{csquotes}
\usepackage{siunitx}
\sisetup{parse-numbers = false}
\DeclareSIUnit{\gal}{Gal}


%\usepackage[style=long]{glossaries}
%\usepackage{abntex2glossaries}
\usepackage{mathtools,bm}
\usepackage{subcaption}
\usepackage{aliascnt}
\usepackage{chngcntr}
\usepackage{listings}
\usepackage{forest}
\usepackage{booktabs}
\usepackage{xparse}
\usepackage{mfirstuc}
% permite representar o cancelamento de termos em texto ou equacoes
\usepackage{cancel}
% cores extendidas
\usepackage{xcolor}
% gera diagramas a partir de listas
\usepackage{smartdiagram}
% Para a figura ficar na posição correta
\usepackage{float}
% supporte para fontes da Text Companion
\usepackage{textcomp}
% uso de longtable
\usepackage{longtable}
% simbolos matematicos
\usepackage{amsmath}
\usepackage{amssymb}
\DeclareMathOperator{\RMSE}{RMSE}
\DeclareMathOperator{\WRMSE}{WRMSE}
% páginas em paisagem
\usepackage{lscape}
% mescla de colunas em tabelas
\usepackage{multicol}
% mescla de linhas em tabelas
\usepackage{multirow}
% criação do indice de quadros
\usepackage{newfloat}
% configura legenda
\usepackage{caption}
%[format=plain]
	%\renewcommand\caption[1]{%
    \captionsetup{font=small}	% tamanho da fonte 10pt
    %,format=hang
 	% \caption{#1}}
	%\captionsetup{width=0.8\textwidth}

  % ----------------------------------------------------------
\usepackage{hyperref}
\hypersetup{
    colorlinks=true,
    linkcolor=blue,     % cor dos links internos (sumário, referências cruzadas etc.)
    citecolor=blue,     % cor das citações (quando se usa biblatex/natbib)
    filecolor=blue,     % cor de links para arquivos locais
    urlcolor=blue       % cor dos hyperlinks (endereço na internet)
}

% ----------------------------------------------------------
%\usepackage{fontspec}
%\newfontfamily\jetbrainsmono{JetBrainsMono NL Mono}
\renewcommand{\ttdefault}{lmvtt}


% ----------------------------------------------------------

% definição das cores
\definecolor{blue}{RGB}{55,10,249}
\definecolor{codegray}{rgb}{0.5,0.5,0.5}
\definecolor{codepurple}{rgb}{0.58,0,0.82}
\definecolor{backcolour}{rgb}{0.97,0.97,0.97}
\definecolor{codeblue}{rgb}{0.26,0.44,0.76}
\definecolor{codegreen}{rgb}{0,0.6,0}
  % ----------------------------------------------------------

\lstset{
    language=Python,
    basicstyle=\fontfamily{lmvtt}\selectfont\small,
    keywordstyle=\color{codeblue}\bfseries,
    stringstyle=\color{codegreen},
    commentstyle=\color{codegray}\itshape,
    numbers=none,
    showstringspaces=false,
    breaklines=true,
    captionpos=t,
    frame=single
}

% Define o novo ambiente flutuante
\DeclareFloatingEnvironment[
  fileext=lon,
  listname={Lista de Funções},
  name={Função },
  placement=tbhp,
]{funcaofloat}

% Faz o \autoref chamar "função" (minusculo, como você queria)
\providecommand*{\funcaofloatautorefname}{função}

\newcommand{\funcauto}[1]{\textbf{função~\ref{#1}}}
\newcommand{\Funcauto}[1]{\textbf{Função~\ref{#1}}}
\newcommand{\claauto}[1]{\textbf{classe~\ref{#1}}}


% Se quiser numeração contínua (sem reset por capítulo)
\counterwithout{funcaofloat}{chapter}


% Pacotes de citações BibLaTeX
% ----------------------------------------------------------
\usepackage[style=authoryear,backref=true,backend=biber,citecounter=true,backrefstyle=three,maxnames=2]{biblatex}

% Remover recuo do sobrenome do autor
\setlength{\bibhang}{0pt}

\DefineBibliographyStrings{brazil}{%
 backrefpage = {Citado \arabic{citecounter} vez na página},% originally "cited on page"
 backrefpages = {Citado \arabic{citecounter} vezes nas páginas},% originally "cited on pages"
}

% ----------------------------------------------------------
\PrepareListOf{quadro}{%
\renewcommand{\cftfigpresnum}{Quadro~}}

\DeclareFloatingEnvironment[
fileext=loq,
listname={\textbf{LISTA DE QUADROS}},
name=Quadro,
%placement=p,
within= none, % numeracao continua
%within=section, % numeracao reinicia em cada seccao
%chapterlistsgaps=off
]{quadro}

\newlistentry{quadro}{loq}{0}

% Define o novo ambiente flutuante
\DeclareFloatingEnvironment[
    fileext=lon,
    listname={Lista de Notebooks},
    name={Jupyter Notebook },  % Aqui define o nome que aparece na legenda
    placement=tbhp,
]{jupyternotebook}

% Define contador auxiliar para integração com autoref
\newaliascnt{jupyternotebookcnt}{jupyternotebook}

% Agora definimos o nome para autoref explicitamente
\providecommand*{\jupyternotebookautorefname}{notebook}

\newcommand{\notebookauto}[1]{\textbf{notebook~\ref{#1}}}
\newcommand{\Notebookauto}[1]{\textbf{Notebook~\ref{#1}}}

\newcommand{\eqauto}[1]{\textbf{equação~\ref{#1}}}

\newcommand{\figauto}[1]{\textbf{figura~\ref{#1}}}
\newcommand{\Figauto}[1]{\textbf{Figura~\ref{#1}}}


% Finaliza o alias (necessário)
\aliascntresetthe{jupyternotebookcnt}

% Remove a numeração por capítulo
\counterwithout{jupyternotebook}{chapter}

% Customize ‘List of Diagrams’
\PrepareListOf{quadro}{%
\renewcommand{\cftquadropresnum}{\normalsize{QUADRO}~}
\setlength{\cftquadronumwidth}{3.2cm}
%\renewcommand{\cftquadroname}{\quadroname\space}
\renewcommand*{\cftquadroaftersnum}{\hfill--\hfill}
}

\makeatletter
%% we define a helper macro for adjusting lists of new floats to
%% accept a * behind them for not being shown in the TOC, like
%% the other list printing commands in memoir
\newcommand{\AdjustForMemoir}[1]{%
  \csletcs{kept@listof#1}{listof#1}%
  \csdef{listof#1}{%
    \@ifstar
     {\csappto{newfloat@listof#1@hook}{\append@star}%
      \csuse{kept@listof#1}}%
     {\csuse{kept@listof#1}}%
  }
}
\def\append@star#1{#1*}
\makeatother


\AdjustForMemoir{quadro} % prepare `\listofdirfigures` so it accepts a *

\makeatletter
\let\oldcontentsline\contentsline
\def\contentsline#1#2{%
    \expandafter\ifx\csname l@#1\endcsname\l@section
    \expandafter\@firstoftwo
    \else
    \expandafter\@secondoftwo
    \fi
    {%
    	\oldcontentsline{#1}{\MakeTextUppercase{#2}}%
    }{%
    \normalsize %ajusta tamanho da fonte na lista
    \oldcontentsline{#1}{#2}%
}%
}
\makeatother


% Ajusta indentação de Referencias no ToC
% ----------------------------------------------------------
\defbibheading{bay}[\bibname]{%
  \chapter*{#1}%
  \markboth{#1}{#1}%
  \addcontentsline{toc}{chapter}
  {\protect\numberline{}\bibname}
}

\makeatletter
\pretocmd{\chapter}{\addtocontents{toc}{\protect\addvspace{-5\p@}}}{}{}
\pretocmd{\section}{\addtocontents{toc}{\protect\addvspace{2\p@}}}{}{}
\makeatother



\newcommand{\DirectoryTree}[1]{%
\begin{forest}
for tree={
    font=\ttfamily,
    grow'=0,
    % padrão para nós não-raiz
    parent anchor=west,
    child anchor=west,
    anchor=west,
    calign=first,
    l=20pt,
    s sep=8pt,
    inner xsep=4pt,
    % desloca o rótulo da raiz e faz com que as arestas dos filhos saiam pela direita
    if level=0{
        xshift=5pt,          % move só o rótulo raiz para a direita
        parent anchor=east,  % faz os filhos se conectarem à direita da raiz
    }{},
    edge path={
        \noexpand\path [draw, \forestoption{edge}]
        (!u.parent anchor) -| + (5pt,0) |- (.child anchor)\forestoption{edge label};
    },
}
#1
\end{forest}
}




% Novo ambiente flutuante para diagramas de diretórios
\DeclareFloatingEnvironment[
    fileext=lod,
    listname={Lista de Sub-diretórios},
    name={Sub--diretório },
    placement=tbhp,
]{directorydiagram}

% Suporte ao \autoref
\newaliascnt{directorydiagramcnt}{directorydiagram}
\providecommand*{\directorydiagramautorefname}{sub-diretório}

\newcommand{\sdauto}[1]{\textbf{sub--diretório~\ref{#1}}}
\newcommand{\Sdauto}[1]{\textbf{Sub--diretório~\ref{#1}}}


\aliascntresetthe{directorydiagramcnt}

% Numeração independente de capítulos
\counterwithout{directorydiagram}{chapter}
